%==============================================================
%  Nwele Uchenna David — Curriculum Vitae
%  Layout/typography replicated from supplied template
%  Compile with: ./build.sh   (or: pdflatex cv.tex)
%==============================================================
\documentclass[11pt,a4paper]{article}

\usepackage[left=0.7in,right=0.7in,top=0.4in,bottom=0.35in]{geometry}
% NOTE: T1 fontenc deliberately omitted - it makes Computer Modern ligatures
% (fi, fl, ff) extract as control characters, garbling words for ATS parsers.
% Default OT1 + glyphtounicode extracts cleanly. Do not re-add fontenc.
\usepackage[utf8]{inputenc}
\usepackage{xcolor}
\usepackage{enumitem}
\usepackage{array}
\usepackage{fancyhdr}
\usepackage[hidelinks]{hyperref}

% - ATS: make PDF text extraction reliable -------------------
\input{glyphtounicode}
\pdfgentounicode=1

\pagestyle{fancy}
\fancyhf{}
\renewcommand{\headrulewidth}{0pt}
\renewcommand{\footrulewidth}{0pt}
\raggedbottom
\setlength{\parindent}{0pt}

% --------------------- EDIT THESE -----------------------------
\newcommand{\myemail}{developeruche@gmail.com}
% Phone: fill in between the braces to show it in the header's second row.
% Left empty, the slot simply collapses and nothing is printed.
\newcommand{\myphone}{}
% --------------------------------------------------------------

\definecolor{sectbg}{gray}{0.86}

% Section header: bordered, gray-filled, sans-serif bold bar
\newcommand{\cvsection}[1]{%
  \vspace{6pt}%
  \noindent\fcolorbox{black}{sectbg}{%
    \parbox[c][19pt][c]{\dimexpr\textwidth-2\fboxsep-2\fboxrule\relax}{%
      \hspace{7pt}\sffamily\bfseries\large #1}%
  }%
  \par\vspace{5pt}%
}

% Outer list of entries (bullet markers)
\newlist{cventries}{itemize}{1}
\setlist[cventries]{label=\textbullet,leftmargin=1.3em,itemsep=3.5pt,%
                   parsep=0pt,topsep=2pt,partopsep=0pt}

% Inner achievement bullets (dash markers)
\newlist{cvbullets}{itemize}{1}
\setlist[cvbullets]{label=\textbf{-},leftmargin=1.4em,itemsep=0.5pt,%
                    parsep=0pt,topsep=3pt,partopsep=0pt}

% Entry heading: {org}{location}{title}{dates}
\newcommand{\cventry}[4]{%
  \begin{tabular*}{\linewidth}[t]{@{\extracolsep{\fill}}l r@{}}
    \textbf{#1} & #2 \\
    \textit{#3} & \textit{#4}
  \end{tabular*}%
}

% Two-column line for open-source / project rows
\newcommand{\cvrow}[2]{%
  \begin{tabular*}{\linewidth}[t]{@{\extracolsep{\fill}}l r@{}}
    \textbf{#1} & \textit{#2}
  \end{tabular*}%
}

\begin{document}

%==============================================================
% HEADER
%==============================================================
\begin{tabular*}{\linewidth}{@{\extracolsep{\fill}}l r@{}}
  {\Huge Nwele Uchenna David} & \myemail \\[3pt]
  Systems Engineer - Rust, Ethereum Execution Layer, zkVMs & \myphone \\[1pt]
  Open to: Remote (global), Relocation
    & \href{https://github.com/developeruche}{github.com/developeruche} \\[1pt]
  \href{https://developeruche.com}{developeruche.com}
    & \href{https://linkedin.com/in/developeruche}{linkedin.com/in/developeruche}
\end{tabular*}

%==============================================================
\cvsection{Work Experience}

\begin{cventries}


%--------------------------------------------------------------
\item \cventry{Ethereum Foundation - zkEVM Team}{Remote}
              {Research Engineer (Contract)}{May - August 2026}
\begin{cvbullets}
  \item Proposed and drafted a REST/SSZ \texttt{newPayloadWithWitness} Engine API endpoint for Reth and
        Ethrex, motivated by EIP-8025 and real-time zkVM attestation; benchmarks showed order-of-magnitude
        witness-delivery latency reduction over the existing round-trip.
  \item Implemented \texttt{zilkworm-stateless}, a C++ zkVM guest program for stateless Ethereum block
        validation, covering EEST spec compliance and Merkle-Patricia Trie changes.
  \item Contributed to \texttt{ere-guests}, a collection of stateless block-validation guest programs for
        Ethereum blocks, and to \texttt{ere}, the Foundation's unified zkVM interface and toolkit.
  \item Integrated a purpose-built TCP engine into Scroll's \texttt{reth-witness-indexer}.
\end{cvbullets}

%--------------------------------------------------------------
\item \cventry{Ethereum Foundation - zkEVM Team}{Remote, Buenos Aires, Argentina}
              {Research Engineer, Ethereum Protocol Fellowship}{June - Nov 2025}
\begin{cvbullets}
  \item Benchmarked precompiles across four major zkVM platforms (RISC0, SP1, OpenVM, Ziren), showing
        that optimized software with lean precompiles outperforms generic implementations with fat
        precompiles by roughly 27$\times$; published the findings as \textit{Fat vs.\ Lean Precompiles in
        zkVMs}, arguing for a precompile strategy that avoids ecosystem fragmentation.
  \item Identified AOT compilation as the primary lever for 10$\times$ zkVM performance gains, reaching
        36M gas in 0.5s on consumer hardware, after profiling trace-generation bottlenecks.
  \item Designed and prototyped a hybrid Ethereum architecture replacing the EVM with a RISC-V execution
        engine while preserving backward compatibility, via magic-number detection and a mini-EVM
        interpreter running as a RISC-V program.
  \item Integrated the hybrid VM into a production Reth node and built \texttt{cargo-hybrid}, a CLI
        toolchain for scaffolding, compiling and deploying RISC-V smart contracts; benchmarked end-to-end
        against \texttt{revm}.
  \item Prototyped ``Stateless Attestors'' with the Prysm team, enabling consensus-layer nodes to verify
        execution payloads through zkEVM proofs without a full execution-layer dependency.
  \item Built a stateless Ethereum block validator running inside \texttt{ziskemu}; compiled Geth
        bare-metal with the Tamago toolchain to eliminate Linux syscall dependencies on zkVM targets.
\end{cvbullets}


%--------------------------------------------------------------
\item \cventry{Independent Research - Rust Systems, zkVMs \& ZKPs}{Lagos, Nigeria}
              {Research Engineer (self-directed)}{June 2024 - June 2025}
\begin{cvbullets}
  \item Stepped back from full-time engineering to work on cryptography and zero-knowledge full time -
        reading papers daily, running research calls with peers, and implementing each result in Rust
        rather than stopping at the write-up; converged on zkVMs as the specialism.
  \item Built the polynomial layer the rest of the work rests on - univariate and multilinear
        representations with their evaluation, interpolation and arithmetic.
  \item Implemented the sum-check protocol and GKR, then succinct GKR - GKR composed with a multilinear
        KZG commitment scheme.
  \item Implemented Groth16, and KZG polynomial commitments in both univariate and multilinear form.
  \item Implemented PLONK end to end - prover, verifier, circuit compiler, and a transpiler lowering
        GKR circuits to Plonkish arithmetization.
  \item Published the whole body of work as \texttt{cryptography-n-zk-research}, an open repository with
        accompanying notes on computational proof systems.
\end{cvbullets}


%--------------------------------------------------------------
\item \cventry{MyAI - Me Protocol}{USA, Remote}
              {Senior Blockchain Engineer}{Nov 2022 - July 2025}
\begin{cvbullets}
  \item Led the blockchain engineering team delivering Me Protocol V3, V4 and V5 alongside the CTO,
        owning architecture through to mainnet deployment.
  \item Architected the protocol in Solidity and Yul on a custom Diamond-standard design, with
        Foundry-based deployment and testing and an in-house migration toolkit supporting both upgrades
        and rollbacks across a system of contracts.
  \item Led development of a custom optimistic rollup in Rust settling to Polygon for asset transfer and
        L2-to-L1 bridging; omitting a sandboxed VM kept it substantially lighter than a general-purpose
        EVM rollup for the workload.
\end{cvbullets}

%--------------------------------------------------------------
\item \cventry{Polytope Labs - Hyperbridge}{Remote}
              {Blockchain Engineer}{Jan - April 2024}
\begin{cvbullets}
  \item Rust engineer on Tesseract, the Hyperbridge cross-chain relayer, and Substrate runtime engineer
        on the Nexus runtime deployed to Polkadot mainnet.
  \item Built HyperClient, a Rust library compiled to WebAssembly that lets Hyperbridge clients manage
        in-flight ISMP requests; implemented teleport support for DOT addresses.
\end{cvbullets}

%--------------------------------------------------------------
\item \cventry{Scroll}{Remote}
              {Smart Contract Engineer (Contract)}{April - June 2024}
\begin{cvbullets}
  \item Built the on-chain contract layer for Scroll Pass, a zero-knowledge event ticketing platform
        delivered with the Ethereum Community Fund.
\end{cvbullets}

%--------------------------------------------------------------
\item \cventry{Nethermind}{Remote}
              {Smart Contract Engineer (Intern)}{Jan - April 2023}
\begin{cvbullets}
  \item \textbf{Valantra:} built a DEX combining a Limit Order Book with constant-product and
        constant-sum AMM algorithms over a shared liquidity base, avoiding fragmentation; exchange, pool
        and management logic fully modularized.
  \item \textbf{Libre:} contributed to a decentralized stock-market protocol for on-chain IPOs.
\end{cvbullets}

%--------------------------------------------------------------
\item \cventry{AfriOne, Nahmii \& Web3Bridge}{Remote / Lagos, Nigeria}
              {Smart Contract Engineer, Design Architect \& Team Lead (Contract)}{2021 - 2023}
\begin{cvbullets}
  \item \textbf{AfriOne:} led the blockchain team; shipped a contract system on a private Geth network
        backing a consumer finance product, plus an SDK for the mobile and backend teams.
  \item \textbf{Nahmii:} built Bondii, an asset-bonding system, and a cross-chain staking platform
        unifying assets across Ethereum L1 and Nahmii 3; ported the Uniswap contracts and UI to Nahmii 3.
  \item \textbf{Web3Bridge:} smart contract team lead for an on-chain governance system built on the Diamond standard.
\end{cvbullets}

\end{cventries}

%==============================================================
\cvsection{Open Source \& Research}

\begin{cventries}

\item \cvrow{Ethereum core tooling in Rust - Alloy, REVM, OpenZeppelin Stylus}{Merged contributions}
\begin{cvbullets}
  \item Contributor to \textbf{Alloy} (the ground-up rewrite of ethers-rs), \textbf{REVM} (the Rust EVM
        used across the client ecosystem), and \textbf{OpenZeppelin's} Stylus contract library.
\end{cvbullets}

\item \cvrow{RAX - RISC-V to x86-64 AOT compiler (Sublinear Labs)}{Contributor}
\begin{cvbullets}
  \item Contributor to \textbf{RAX}, a RISC-V (RV64IMA) to x86-64 ahead-of-time compiler for ELF
        binaries with interpreter and JIT backends; contributed Ethereum block-execution guest
        programs, stateless block-execution handling, and the test/benchmark flow.
\end{cvbullets}

\item \cvrow{RISC-V EVM Experiment - research \& Rust proof of concept}{2025}
\begin{cvbullets}
  \item Evaluated RISC-V as an alternative execution environment for smart contracts; published
        \textit{Bridging Worlds: A Performance and Feasibility Analysis of RISC-V Integration with the
        Ethereum Virtual Machine}. Shipped a supporting RV32IM assembler and emulator in Rust.
\end{cvbullets}

\end{cventries}

%==============================================================
\cvsection{Skills}

\begin{cventries}[itemsep=3pt]
  \item \textbf{Languages:} Rust, Solidity, C++, TypeScript, Go
  \item \textbf{Ethereum:} Reth, REVM, Alloy, Foundry, Engine API, SSZ, Merkle-Patricia Tries, EIP authorship
  \item \textbf{Zero-Knowledge:} zkVMs (RISC0, SP1, OpenVM, Ziren, ZisK), Groth16, PLONK, KZG, GKR, sum-check
  \item \textbf{Systems:} RISC-V (RV32IM), compilers, emulators, AOT compilation, performance benchmarking, WebAssembly
  \item \textbf{Distributed Systems:} Substrate / Polkadot SDK, optimistic rollups, cross-chain messaging (ISMP)
\end{cventries}

%==============================================================
\cvsection{Education}

\begin{cventries}
  \item \cvrow{University of Benin - B.Eng, Computer Engineering}{Benin City, Nigeria \quad 2019 - 2024}
  \item \cventry{Polkadot Blockchain Academy}{Singapore}
                {Blockchain Development - Rust, Substrate Runtime Engineering}{May - June 2024}
\end{cventries}

\end{document}
